\section{The initial family of regions}
\label{sec:area_law_initial_regions}

Fix an arbitrary origin $o\in\mathbb R^2$, a finite endpoint set
$Z\subset\mathbb Z^2$, an integer width $C\ge 2$ and an initial layer
$k_0\ge 50{,}000{,}000$. At each layer $k\ge k_0$, choose arbitrary
residues $a_k,b_k$ modulo $2^{p_k-\ell_k}$. Retain the actual primary
indices, belt cells, actual midpoint subdivisions and triangle colors of
\cite[Section~11]{OpenAI2026AreaLaw}.

% Source: 10-geometry.tex, prop:two-families, lines 154--177,
% 212--218 and 299--323, revision adc7f1241b42e322a6451854ab7e4b4c146bf78a.
% These entries describe the initial family before local and recursive repairs.
% Closed coverage alone does not give disjointness or lattice open coverage.

\begin{definition}[Initial region identifiers]
    \label{def:area_law_initial_region_indices}
    \lean{TNLean.PEPS.AreaLaw.Geometry.InitialRegionIndex}
    \leanok
    \uses{def:area_law_primary_pitch_indices,def:area_law_belt_fan_color,
        def:area_law_fan_runs}
    Let $G_k$ be the retained primary pitch indices and $B_k$ the actual
    belt-cell indices at layer $k$. For $z\in B_k$, let $\mathcal R_{k,z}$
    be the runs of its actual colored fan. The identifier family is the
    disjoint union
    \begin{align}
        \mathcal I
         & =\{*\}\sqcup\coprod_{k\ge k_0}G_k
        \sqcup\coprod_{k\ge k_0}\coprod_{z\in B_k}\mathcal R_{k,z}.
    \end{align}
    The symbol $*$ identifies the dummy region. A primary identifier
    retains its layer and pitch index, including when its region has
    disconnected pieces. A run identifier retains its layer and cell.
    Distinct runs remain distinct even when their colors agree.
\end{definition}

\begin{definition}[Closed birth regions]
    \label{def:area_law_initial_birth_regions}
    \lean{TNLean.PEPS.AreaLaw.Geometry.initialBirthRegion}
    \leanok
    \uses{def:area_law_initial_region_indices,def:area_law_dyadic_layers,
        def:area_law_primary_regions,def:area_law_fan_runs}
    Write $N_0=N_{k_0}$ for the initial half-open neighborhood,
    $R_{k,J}$ for a primary birth region and $T_{k,z,R}$ for a belt-run
    region. For $i\in\mathcal I$, define
    \begin{align}
        X_i^0 & =
        \begin{cases}
            \overline{N_0}, & i=*,                                         \\
            R_{k,J},        & i=(k,J),\ J\in G_k,                          \\
            T_{k,z,R},      & i=(k,z,R),\ z\in B_k,\ R\in\mathcal R_{k,z}.
        \end{cases}
    \end{align}
\end{definition}

\begin{definition}[Initial open interiors]
    \label{def:area_law_initial_open_regions}
    \lean{TNLean.PEPS.AreaLaw.Geometry.initialOpenRegion}
    \leanok
    \uses{def:area_law_initial_birth_regions}
    The initial open region with identifier $i$ is
    $\operatorname{int}(X_i^0)$.
\end{definition}

\begin{definition}[Initial region colors]
    \label{def:area_law_initial_region_colors}
    \lean{TNLean.PEPS.AreaLaw.Geometry.initialRegionColor}
    \leanok
    \uses{def:area_law_initial_region_indices,def:area_law_belt_fan_color,
        thm:area_law_fan_run_color}
    Assign $\lambda(i)\in\mathbb Z/2\mathbb Z$ by
    \begin{align}
        \lambda(i) & =
        \begin{cases}
            k_0-1,                                 & i=*,       \\
            k,                                     & i=(k,J),   \\
            \text{the common triangle color in }R, & i=(k,z,R).
        \end{cases}
    \end{align}
    Run-color constancy makes the last value well defined.
\end{definition}

\begin{theorem}[Agreement with every constituent triangle]
    \label{thm:area_law_initial_run_color}
    \lean{TNLean.PEPS.AreaLaw.Geometry.initialRegionColor_beltRun_eq}
    \leanok
    \uses{def:area_law_initial_region_colors}
    If $e$ is any constituent triangle of an actual belt run $R$ in
    layer $k$ and cell $z$, then $\lambda(k,z,R)$ equals the actual
    color of $e$.
\end{theorem}
\begin{proof}\leanok
    \uses{thm:area_law_fan_run_color,def:area_law_initial_region_colors}
    The run is the component containing $e$. Its common color is therefore
    the color of $e$, by constancy on that component.
\end{proof}

\begin{theorem}[Closed cover by the initial family]
    \label{thm:area_law_initial_closed_cover}
    \lean{TNLean.PEPS.AreaLaw.Geometry.initialBirthRegions_cover}
    \leanok
    \uses{def:area_law_initial_birth_regions}
    If $Z\ne\varnothing$, then the closed birth regions cover the plane:
    \begin{align}
        \bigcup_{i\in\mathcal I}X_i^0 & =\mathbb R^2.
    \end{align}
\end{theorem}
\begin{proof}\leanok
    \uses{thm:area_law_unique_fine_cell,
        thm:area_law_nonbelt_unique_primary,thm:area_law_fan_run_cover,
        def:area_law_initial_birth_regions}
    Fix $x\in\mathbb R^2$. If $x\in N_0$, then
    $x\in\overline{N_0}=X_*^0$. Otherwise fine-cell exhaustion gives
    an actual fine cell at some layer $k\ge k_0$ containing $x$.
    Its closure also contains $x$. If the cell is in the belt, its
    closed square is covered by its run regions, so $x$ lies in one
    $T_{k,z,R}$. If the cell is outside the belt, the unique retained
    primary containing its closure also contains $x$. Each case places
    $x$ in a birth region with an actual identifier.
\end{proof}
